\chapter{Night Sweats}

\section*{Definition}
Nocturnal sweating refers to excessive perspiration during sleep saturating bedclothes, representing abnormal thermoregulatory response indicating underlying physiological or pathological processes.

\section*{What Happens in Your Body}
Night sweats result from hypothalamic thermoregulatory dysregulation during sleep when autonomic control predominates. Increased sympathetic outflow triggers eccrine gland activation. Malignancy-related sweats involve immune signaling proteins effects (TNF-$alpha$, IL-1) altering set-point. Hormonal fluctuations decrease thermoneutral zone causing vasomotor instability.

\section*{Causes}
\begin{itemize}
  \item Lymphoma (Pel-Ebstein syndrome alternating fevers/night sweats)
  \item Tuberculosis (evening fevers and night sweats)
  \item HIV/AIDS
  \item Hyperthyroidism (Graves disease)
  \item Menopause/vasomotor symptoms
  \item Antidepressants (SSRIs, SNRIs, TCAs)
\end{itemize}

\section*{When to See a Doctor}
See your doctor if:
\begin{itemize}
  \item Drenching night sweats requiring bedclothes change
  \item Unexplained weight loss with night sweats
  \item Fever accompanying night sweats
  \item Night sweats in malignancy patient
  \item Night sweats with palpitations/tremor
\end{itemize}

\section*{Related Symptoms}
\begin{itemize}
  \item Weight loss
  \item Fever
  \item Fatigue
  \item Malaise
  \item Lymphadenopathy
  \item Hot flashes
  \item Anxiety
  \item Cough
\end{itemize}
\textbf{Important:} Night sweats with unexplained weight loss and fever constitute red flags for lymphoma or tuberculosis. Investigations include CBC, ESR/CRP, thyroid function, HIV test, chest radiograph.

\section*{Self-Care / Home Management}
\begin{itemize}
  \item Rest and maintain adequate hydration with water, clear fluids, or oral rehydration solutions.
  \item Use acetaminophen or ibuprofen to reduce fever and body aches (follow label dosing instructions).
  \item Dress lightly and keep the room at a comfortable temperature — avoid bundling excessively.
  \item Monitor temperature at regular intervals and track other symptoms to identify worsening trends.
  \item Seek medical attention for fever above 40	extdegree{}C (104	extdegree{}F), fever lasting more than 3 days, or fever with stiff neck, severe headache, rash, or difficulty breathing.
\end{itemize}

\section*{How Your Doctor Will Evaluate You}
\begin{itemize}
  \item History includes fever duration, pattern (intermittent, remittent, continuous), associated symptoms (chills, rigors, sweats), and exposure history.
  \item Vital signs assess fever severity and hemodynamic stability; complete physical examination seeks a source.
  \item Laboratory evaluation includes CBC with differential, blood cultures, urinalysis, inflammatory markers (CRP, ESR), and targeted testing based on exposure history.
  \item Imaging (chest X-ray, abdominal ultrasound, CT) is guided by clinical findings and suspected source.
\end{itemize}

\section*{Treatment Options}
\begin{itemize}
  \item Fever-reducing medications (acetaminophen, NSAIDs) provide symptomatic relief but do not treat the underlying cause.
  \item Antibiotics based on the most likely cause may be indicated for suspected bacterial infection after appropriate cultures are obtained.
  \item Antiviral, antifungal, or antimalarial therapy is directed at specific pathogens based on exposure and testing.
  \item Source control (drainage of abscess, removal of infected device) is essential for focal infections.
\end{itemize}

\section*{References}
\begin{thebibliography}{99}
\bibitem{ref1} Viera AJ, Bond MM, Yates SW. "Diagnosing night sweats." \textit{American Family Physician}, 2003;67(5):1019--1024.
\bibitem{ref2} Mold JW, Lawler F, Roberts M. "The prevalence and clinical significance of night sweats in a general practice population." \textit{Journal of Family Practice}, 2006;55(8):694--699.
\bibitem{ref3} Siebert J, Siebert M. "Night sweats: a systematic review." \textit{Journal of the American Board of Family Medicine}, 2011;24(4):434--442.
\bibitem{ref4} Wilborn SL, Gupta S. "Night sweats: a guide for the primary care clinician." \textit{Clinical Medicine Insights: Therapeutics}, 2015;7:23--30.
\bibitem{ref5} Mikkelsen LH, Pedersen LM. "Night sweats in the cancer patient." \textit{Supportive Care in Cancer}, 2021;29(4):1803--1810.
\bibitem{ref6} Freedman RR. "Physiology of hot flashes." \textit{American Journal of Human Biology}, 2001;13(4):453--464.

\end{thebibliography}

\clearpage
